\chapter{If statements}

\section{If condition}

If statements are weird in bash. First of all, bash is very space sensitive so manage with care.
Then the syntax is very atypical.

\lstinputlisting[
    language=bash,
    caption={4-bash.sh},
    captionpos=t,
    label={lst:4-bash},
    breaklines=true]
{../4-bash/4-bash.sh}

Now, you might be wondering why are we using a command flag inside our if statement? And the reason
for that is because actually, the square bracket is a command itself. It's an abreviation of the
test command. The -eq flag stands for equals to and expects an argument afterward. The closing
square bracket is just expected by the alias definition of the test command.

\section{A more modern approach for number comparisson}

Now, if you wanted to use your typical ==, !=, <, etc operators,you would use a double parentheses
command. This is as modern as the eyar 1997 can be labeled modern, which is not at all. So use this
if you want a more normal experience.

\lstinputlisting[
    language=bash,
    caption={4-bash.sh},
    captionpos=t,
    label={lst:4.1-bash},
    breaklines=true]
{../4-bash/4.1-bash.sh}
\pagebreak
\section{Else statement}

\lstinputlisting[
    language=bash,
    caption={4-bash.sh},
    captionpos=t,
    label={lst:4.2-bash},
    breaklines=true]
{../4-bash/4.2-bash.sh}

\section{Not operator}

\lstinputlisting[
    language=bash,
    caption={4-bash.sh},
    captionpos=t,
    label={lst:4.3-bash},
    breaklines=true]
{../4-bash/4.3-bash.sh}

\section{-f flag}

Since the `[]' operator is a bash command, it can interact with directories, files and such.
This means we can compare agains files or directories directly. The `-f' flag is just one way
to do this, where it checks wether a path exists AND is a file. if you just wanted to see if 
it existed as a file or directory, you'd use the -e flag.

\lstinputlisting[
    language=bash,
    caption={4-bash.sh},
    captionpos=t,
    label={lst:4.4-bash},
    breaklines=true]
{../4-bash/4.4-bash.sh}
\pagebreak
\section{A more ``practical'' example}

We want to check if a command exists, if it does run it, otherwise install it and then run it.

\lstinputlisting[
    language=bash,
    caption={4-bash.sh},
    captionpos=t,
    label={lst:4.5-bash},
    breaklines=true]
{../4-bash/4.5-bash.sh}

\section{The better way of doing it}

We used the test command to check wether a file existed in our system. However, this is a very
general way of doing it, the shell has built in ways of managin this specific situation.
We want to check not if a file exist but if a command exists. So we'll use a Shell builtin
known as a command. Shell builtins are like any other commands but they are baked directly into
the shell as oposed of being binary files. A common example of this would be the cd command.
If you wish to know whether a command is a shell builtin, use the type command.

\lstinputlisting[
    language=bash,
    caption={4-bash.sh},
    captionpos=t,
    label={lst:4.6-bash},
    breaklines=true]
{../4-bash/4.6-bash.sh}

The reason why you wouldn't want to use which is because the which command doesn't contemplate 
shell buitlin commands so it might not work with every command. Also command tends to be faster 
since it doesn't require to open up a file and do many processes but is baked into the running process.
known as the shell.
\pagebreak

\section{Logic operators}

We already saw the AND operator. Naturally, it can also be used inside the if statement argument.

\lstinputlisting[
    language=bash,
    caption={4-bash.sh},
    captionpos=t,
    label={lst:4.7-bash},
    breaklines=true]
{../4-bash/4.7-bash.sh}

\section{Else if}

Of course any well self respecting language would have the inclussion of the humble else if.

\lstinputlisting[
    language=bash,
    caption={4-bash.sh},
    captionpos=t,
    label={lst:4.8-bash},
    breaklines=true]
{../4-bash/4.8-bash.sh}


\section{Ok but how does the if statement actually work?}

The if statements in bash scripting don't work with your typical boolean values, they are not expecting
a true nor a false. What they work with is exit codes. If a command ran successfully, it will return a 0 exit code
otherwise, any other number. So 0 behaves like true, any other value behaves like false. This is the 
opposite of how you'd expect it to work. More on exit codes next chapter.